\documentclass[11pt]{article}
\usepackage[a4paper,margin=1.06in]{geometry}
\usepackage[T1]{fontenc}
\usepackage[utf8]{inputenc}
\usepackage{lmodern}
\usepackage{amsmath,amssymb,amsthm,mathtools}
\usepackage{enumitem}
\usepackage[hidelinks]{hyperref}
\usepackage{microtype}
\setlength{\emergencystretch}{2em}

\newtheorem{theorem}{Theorem}[section]
\newtheorem{proposition}[theorem]{Proposition}
\newtheorem{lemma}[theorem]{Lemma}
\newtheorem{corollary}[theorem]{Corollary}
\newtheorem{definition}[theorem]{Definition}
\newtheorem{remark}[theorem]{Remark}
\newtheorem{example}[theorem]{Example}

\DeclareMathOperator{\End}{End}
\DeclareMathOperator{\rad}{rad}
\DeclareMathOperator{\Spec}{Spec}
\DeclareMathOperator{\Span}{span}

\title{Presentation-Langlands I: Proof-Carrying Hecke Records}
\author{Luca Blanchi}
\date{June 2026}

\begin{document}
\maketitle

\begin{abstract}
This note develops a finite presentation-theoretic format for Hecke eigensystem data.  In a finite Hecke laboratory, a computed eigensystem should not be recorded only as a list of eigenvalues.  It should also carry a declared window, a separating set of Hecke observables, optional synthetic fingerprints, projector witnesses, negative ideal-membership witnesses for absent candidates, residual fibres, and normalization conventions.

The main algebraic facts are elementary but useful.  Selecting a minimum natural Hecke fingerprint is a weighted set-cover problem on pairs of eigensystems.  If a finite set of commuting semisimple Hecke operators separates the eigensystems in a window, a generic linear combination has simple spectrum on that window.  Spectral interpolation then gives idempotent projectors, while absent eigenvalue tuples are witnessed by finite ideal-membership identities in the Hecke algebra.  In non-semisimple or integral settings, ordinary eigenvalues see only the semisimple quotient; nilpotent and congruence data must be carried separately as residual structure.

The result is a proof-carrying record schema for finite Hecke data.  It is intended as the first component of a broader Presentation-Langlands layer: a way of saying exactly what a computation observes, what it verifies, and what ambiguity remains.
\end{abstract}

\section{Introduction}

Computations with modular forms, automorphic forms, and arithmetic databases frequently produce finite lists of Hecke eigenvalues.  Such lists are often used as identifiers.  The presentation-theoretic question is more precise.  Which finite data identify the object in the declared window?  What verifies that identification claim?

This note answers that question in the simplest setting: a finite-dimensional vector space carrying a finite commutative algebra of computed Hecke operators.  The point is not to reprove multiplicity one or the theory of Hecke algebras.  The point is to package a finite computation so that its identification claim is auditable.

A mature Hecke record should contain:
\begin{enumerate}[label=\textup{(\roman*)}]
\item the finite window in which the claim is made;
\item the Hecke observables used to distinguish the objects;
\item the cost or budget attached to those observables;
\item positive verification data, such as projectors or interpolation identities;
\item negative verification data excluding absent candidates;
\item a residual ledger describing equivalences or unresolved fibres;
\item normalization conventions.
\end{enumerate}

This is a Presentation-Langlands viewpoint: Langlands data are treated not only as objects and correspondences, but as presented, observed, verified, and transferred data.

\section{Finite Hecke laboratories}

\begin{definition}[Finite Hecke laboratory]
A finite Hecke laboratory is a tuple
\[
\mathcal H=(M,A,\mathcal T,\mathcal W,\sim,\mathsf N),
\]
where \(M\) is a finite-dimensional vector space over a field \(k\), \(A\subseteq \End_k(M)\) is a finite commutative \(k\)-algebra generated by a finite list
\[
\mathcal T=\{T_1,\ldots,T_m\},
\]
\(\mathcal W\) is the finite window of eigensystems or generalized eigensystems under consideration, \(\sim\) is the admissible equivalence relation in the window, and \(\mathsf N\) records normalization conventions.
\end{definition}

Typical equivalence relations include equality, Galois conjugacy, unramified twist, packet equivalence, or equality after a declared quotient of the data.  The normalization ledger \(\mathsf N\) records choices such as Hecke operator conventions, local reciprocity convention, measures, local factors, and coefficient-field embeddings.

The coefficient field matters.  The laboratory may be defined over a coefficient field \(k\), over a splitting field \(K/k\), over an algebraic closure, or over an integral coefficient ring \(R\) with fraction field \(K\).  The finite linear-algebra statements below are always made after declaring the field in which eigenvalues are compared.  Integral, congruence, and generalized-eigenspace data are recorded separately from ordinary eigenvalues.

For a finite subset \(S\subseteq \mathcal T\), define the observation map
\[
\mathcal O_S:\mathcal W\longrightarrow \prod_{T\in S} k,\qquad
\lambda\longmapsto(\lambda(T))_{T\in S}.
\]
The residual fibre of \(\lambda\) is
\[
F_S(\lambda)=
\{\mu\in\mathcal W:\mathcal O_S(\mu)=\mathcal O_S(\lambda)\}.
\]
The residual ledger is
\[
\mathsf{Led}_S(\lambda)=F_S(\lambda)/\sim .
\]

An identification claim is complete in the window only when
\[
|\mathsf{Led}_S(\lambda)|=1.
\]
When this fails, the ledger is not an error; it is the mathematically honest statement of what the chosen data do not see.

\section{Natural fingerprints}

Assume first that \(A\) is semisimple after scalar extension to a declared splitting field \(K\), and that the window consists of distinct eigensystems
\[
\lambda_1,\ldots,\lambda_r.
\]
For each operator \(T_i\), define the set of pairs it separates:
\[
C_i=\{\{a,b\}:1\leq a<b\leq r,\ \lambda_a(T_i)\neq\lambda_b(T_i)\}.
\]
Let
\[
U=\{\{a,b\}:1\leq a<b\leq r\}.
\]

\begin{proposition}[Hecke fingerprint as set cover]
A subset \(J\subseteq\{1,\ldots,m\}\) separates the eigensystems in the window if and only if
\[
\bigcup_{i\in J} C_i=U.
\]
Consequently the problem of choosing a minimum-cost natural Hecke fingerprint is exactly a weighted set-cover problem on the pair set \(U\).
\end{proposition}

\begin{proof}
The operators indexed by \(J\) separate the eigensystems precisely when, for every pair \(\lambda_a\neq\lambda_b\), at least one \(T_i\) with \(i\in J\) has different eigenvalues on the two eigensystems.  This is exactly the condition that every pair \(\{a,b\}\) belongs to some \(C_i\) with \(i\in J\).  Adding weights \(c(T_i)\) gives the weighted set-cover formulation.
\end{proof}

This proposition is elementary, but it is operational.  It turns the vague instruction ``compute enough Hecke operators'' into a finite optimization problem with a declared universe of pairs.

\section{Synthetic fingerprints and projectors}

Natural fingerprints use selected Hecke operators.  Sometimes it is convenient to compress them into one synthetic operator.

\begin{proposition}[Generic simple fingerprint]
Let \(T_1,\ldots,T_m\) be commuting semisimple operators whose joint eigenvalue tuples distinguish
\(\lambda_1,\ldots,\lambda_r\) over a declared splitting field.  If the coefficient field is infinite, then a Zariski-generic linear combination
\[
T_{\mathrm{fp}}=\sum_{i=1}^m c_iT_i
\]
has distinct eigenvalues on \(\lambda_1,\ldots,\lambda_r\).  If the coefficient field is finite, the same statement holds after extending scalars to a sufficiently large finite extension, or after choosing the coefficients \(c_i\) in an infinite coefficient field containing the eigenvalues.
\end{proposition}

\begin{proof}
For a pair \(a\neq b\), the equality of the synthetic eigenvalues is
\[
\sum_i c_i\lambda_a(T_i)=\sum_i c_i\lambda_b(T_i).
\]
Since the joint eigenvalue tuples are distinct, this is a proper hyperplane in the coefficient space.  There are finitely many pairs, hence finitely many such hyperplanes.  Any coefficient vector outside their union gives distinct synthetic eigenvalues.
\end{proof}

Suppose \(T_{\mathrm{fp}}\) has distinct eigenvalues
\[
\theta_j=\lambda_j(T_{\mathrm{fp}}).
\]
Define the interpolation polynomial
\[
P_j(x)=
\prod_{i\neq j}\frac{x-\theta_i}{\theta_j-\theta_i}.
\]

\begin{proposition}[Projector witnesses]
In the semisimple window, over a coefficient field in which all differences \(\theta_j-\theta_i\) are invertible,
\[
P_j(T_{\mathrm{fp}})
\]
is the projector onto the \(\lambda_j\)-eigenspace.
\end{proposition}

\begin{proof}
The polynomial \(P_j\) satisfies \(P_j(\theta_j)=1\) and \(P_j(\theta_i)=0\) for \(i\neq j\).  Since \(T_{\mathrm{fp}}\) acts diagonally on the eigenspace decomposition, \(P_j(T_{\mathrm{fp}})\) is the identity on the \(\lambda_j\)-summand and zero on the other summands.
\end{proof}

Thus a proof-carrying record can store not only the values of a fingerprint, but an idempotent witness:
\[
P_j(T_{\mathrm{fp}})^2=P_j(T_{\mathrm{fp}}).
\]

\section{Negative witnesses}

A finite Hecke record should also be able to exclude candidates.
Let \(\alpha=(\alpha_1,\ldots,\alpha_m)\) be a proposed eigensystem on the generators.  Consider the ideal
\[
I_\alpha=(T_1-\alpha_1,\ldots,T_m-\alpha_m)\subseteq A.
\]

\begin{proposition}[Finite Hecke Nullstellensatz witness]
Assume \(A\) is a finite-dimensional commutative semisimple algebra over an algebraically closed field.  The candidate \(\alpha\) is absent from \(\Spec A\) if and only if
\[
1\in I_\alpha.
\]
Equivalently, there exist polynomials or algebra elements \(Q_i\in A\) such that
\[
1=\sum_i Q_i(T_i-\alpha_i).
\]
\end{proposition}

\begin{proof}
In a finite reduced algebra over an algebraically closed field, maximal ideals correspond to characters.  The tuple \(\alpha\) occurs precisely when the ideal generated by the \(T_i-\alpha_i\) is contained in a maximal ideal.  If no such character occurs, the ideal is contained in no maximal ideal, hence it is the whole algebra.  This is equivalent to the displayed identity.
\end{proof}

The displayed identity is finite verification data.  It can be checked by multiplication in \(A\), or by matrices representing the \(T_i\).  It is therefore a natural negative witness for a proof-carrying Hecke record.

\section{Non-semisimple residual data}

Many arithmetic situations are not semisimple over the coefficient ring of interest.  Torsion, congruences, generalized eigenspaces, and integral structures can carry nilpotent information invisible to ordinary eigenvalues.

Let \(A\) be a finite commutative algebra and let
\[
\mathfrak r=\rad(A).
\]
The semisimple quotient is
\[
A_{\mathrm{ss}}=A/\mathfrak r.
\]
Ordinary eigenvalue data factor through \(A_{\mathrm{ss}}\).

\begin{proposition}[Semisimple quotient and nilpotent ledger]
Every eigenvalue record for \(A\) induces an eigenvalue record for \(A_{\mathrm{ss}}\).  The extra information required to recover the full finite algebra from its semisimple quotient is residual nilpotent data, such as the radical filtration
\[
\mathfrak r\supseteq \mathfrak r^2\supseteq\cdots .
\]
\end{proposition}

\begin{proof}
Every character \(A\to \bar k\) vanishes on the nilradical and hence factors through \(A_{\mathrm{ss}}\).  Therefore ordinary eigenvalues depend only on the semisimple quotient.  Any distinction between two algebras, modules, or generalized eigenspaces with the same semisimple quotient must be recorded by data not seen by characters, for instance radical filtrations, generalized eigenspaces, Fitting ideals, extension groups, or other derived observables.
\end{proof}

This is a residual ledger, not a defect of the method.  It says exactly what the chosen observables can and cannot prove.

\subsection{Limits of semisimple records}

A semisimple eigenvalue record proves only a statement about the semisimple quotient and the declared finite window.  Equal ordinary eigenvalues do not by themselves identify integral structures, congruence modules, radical filtrations, or generalized eigenspaces.  When such data matter, the record must include additional entries such as Fitting ideals, radical-layer dimensions, congruence modules, extension data, or generalized-eigenspace witnesses.

\section{Proof-carrying Hecke records}

\begin{definition}[Proof-carrying Hecke record]
A proof-carrying Hecke record for an eigensystem \(\lambda\) in a finite laboratory is a tuple
\[
\mathsf{HeckeRec}(\lambda)=
(M,A,\mathcal W,S,T_{\mathrm{fp}},P_\lambda,\mathsf W^-,
\mathsf{Led},\mathsf N,D_{\mathrm{form}}).
\]
Here \(S\) is a natural fingerprint, \(T_{\mathrm{fp}}\) is an optional synthetic fingerprint, \(P_\lambda\) is a positive projector witness when available, \(\mathsf W^-\) is a family of negative witnesses excluding declared absent candidates, \(\mathsf{Led}\) is the residual ledger, \(\mathsf N\) records normalizations, and \(D_{\mathrm{form}}\) records the expected level of machine or formal verification.
\end{definition}

The record is valid in the declared window if the following checks pass:
\begin{enumerate}[label=\textup{(\roman*)}]
\item the window \(\mathcal W\) and equivalence relation \(\sim\) are declared;
\item the fingerprint \(S\) has the stated residual fibres;
\item projectors satisfy the required interpolation identities;
\item negative witnesses satisfy their ideal-membership identities;
\item residual nilpotent, twist, packet, or Galois-orbit ledgers are explicitly stated;
\item all normalizations used to compare values are recorded.
\end{enumerate}

\begin{theorem}[Finite Hecke record theorem]
In every finite semisimple Hecke laboratory whose computed Hecke observables separate the declared window, each eigensystem admits a proof-carrying Hecke record with a natural fingerprint, a synthetic fingerprint after a generic linear combination over a suitable coefficient field, projector witnesses, negative ideal-membership witnesses for the declared finite list of absent candidates, and a trivial residual ledger modulo the declared equivalence.
\end{theorem}

\begin{proof}
The natural fingerprint is obtained by choosing a set cover on the pair set of the window.  The generic synthetic fingerprint and the projectors are given by the preceding propositions.  Negative witnesses are supplied by the finite Hecke Nullstellensatz proposition for absent candidates.  Since the chosen observables separate the window modulo the declared equivalence, the residual ledger has one class for each identified object.
\end{proof}

\section{Presentation-theoretic role}

The construction isolates a reusable pattern:
\[
\text{finite window}
\quad\longrightarrow\quad
\text{observables}
\quad\longrightarrow\quad
\text{residual fibre}
\quad\longrightarrow\quad
\text{verification data}.
\]
It is not a replacement for automorphic theory.  It is a disciplined layer on top of finite computations and finite consequences of theory.  It makes explicit when a list of eigenvalues identifies an object, when it identifies only a quotient, and what additional data would be needed to resolve the remaining fibre.

This record format is the starting point for the next applications: class-group tomography, transfer audits, robust analytic verification, and proof-carrying links in arithmetic databases.

\begin{thebibliography}{99}

\bibitem{cremona2026}
J. E. Cremona,
\emph{Hecke operators, Hecke eigensystems, and formal modular symbols over number fields},
arXiv:2601.17524, 2026.

\bibitem{lmfdb}
J. Cremona, J. Balakrishnan, N. Dunfield, M. Harrison, D. Roberts, W. Stein, and others,
\emph{The L-functions and modular forms database project},
arXiv:1511.04289, 2015.

\bibitem{diamond-shurman}
F. Diamond and J. Shurman,
\emph{A First Course in Modular Forms},
Graduate Texts in Mathematics 228, Springer, 2005.

\bibitem{miyake}
T. Miyake,
\emph{Modular Forms},
Springer Monographs in Mathematics, Springer, 2006.

\end{thebibliography}

\end{document}
